\section{Discussion critique : biais, limites et plafond de performance}
\label{sec:biais}

L'évaluation de la performance d'un agent ne peut être dissociée des artefacts de son banc de test. L'observation d'un plafond de succès autour de 33\,\% à 37\,\% pour \texttt{Qwen3-14B}, même sous budget de raisonnement, s'explique par la conjonction de facteurs intrinsèques au modèle et de biais inhérents au benchmark.

\subsection{Analyse du plafond empirique : pourquoi le gain de raisonnement reste modéré \texorpdfstring{(+2\,\%)}{(+2\%)}}

Alors que l'activation du raisonnement produit une hausse sensible de consistance ($\passk{2}$, concordance DB), le taux de réussite global stagne sous la barre des 40\,\%. Trois facteurs explicatifs majeurs se dégagent :
\begin{enumerate}
  \item \textbf{Capacité intrinsèque et taille de paramètres (14B face aux modèles frontières)} : la politique \texttt{airline} compte plus de 1\,300 mots de règles conditionnelles denses (barèmes de pénalités selon 4 classes tarifaires, conditions d'annulation proratisées, règles de transfert entre réservations distinctes). Un modèle de 14 milliards de paramètres présente des limites dans la déduction logique croisée et le calcul arithmétique multi-variables sous longue fenêtre de contexte (32k), comparativement à des modèles de 70B+ ou des modèles propriétaires frontières.
  \item \textbf{Impact de la quantification 4-bit (\texttt{Q4\_K\_M})} : la compression à 4 bits induit une dégradation fine sur les calculs numériques précis et la sérialisation exacte des arguments d'outils, ce qui amplifie les échecs lors des calculs de frais de dossier.
  \item \textbf{Rôle et bruit de l'utilisateur simulé} : dans une évaluation de bout en bout, l'agent dépend entièrement de la clarté des réponses du simulateur utilisateur. Lorsque l'utilisateur omet de fournir un identifiant requis, contredit son propre profil ou abrège la session, l'échec est imputé à l'agent bien qu'il résulte d'un dysfonctionnement du simulateur.
\end{enumerate}

\subsection{Biais structurels du système d'évaluation}

\paragraph{Un score centré exclusivement sur l'état final.} 
La composante \texttt{DB} ignore totalement la qualité et l'éthique de la conversation. Deux agents aux comportements diamétralement opposés, l'un refusant fermement une demande illégitime en expliquant la règle, l'autre acquiesçant servilement mais sans déclencher d'outil, reçoivent la même récompense $r = 1$ si la base reste intacte (section~\ref{sec:prise-en-main}).

\paragraph{Le dogme de la trajectoire unique et la pénalité de sur-génération.}
En comparant les états finaux par hachage cryptographique strict, le benchmark postule qu'il n'existe qu'un unique état valide. Dès lors qu'un agent enrichit un appel avec des arguments optionnels valides (tâche \texttt{mock}), il est comptabilisé comme défaillant. Paradoxalement, un modèle doté d'un budget de raisonnement poussé, enclin à contextualiser ses actions, risque d'être pénalisé par cette rigidité formelle.

\paragraph{Fragilité de la métrique agrégée.}
Le score global masque des disparités criantes entre domaines : \texttt{airline} mesure l'état final et des sous-chaînes, \texttt{retail} sollicite un juge LLM pour arbitrer le dialogue, tandis que \texttt{telecom} applique des prédicats d'environnement combinés à des listes d'actions obligatoires. Fusionner ces grandeurs dans un score unifié n'a pas de cohérence statistique.

\subsection{Validité statistique et reproductibilité}

Avec un jeu de 50 tâches, l'intervalle de confiance à 95\,\% pour un taux de succès de $\sim 30\,\%$ est d'environ $\pm 12{,}7$ points sur un essai unique. L'utilisation de $n = 3$ essais et de la métrique $\passk{k}$ est indispensable pour stabiliser l'estimation : nos résultats expérimentaux montrent que l'effet du raisonnement se matérialise principalement sur la fiabilisation des succès conjoints ($\passk{2}$ passant de 18{,}0\,\% à 21{,}3\,\%), ce qu'un test classique à essai unique ($n=1$) n'aurait pas permis de discriminer avec certitude.
